\section{微商概念及其计算}\label{sec:4.1}

\subsection{微商概念}

微商是数学分析最重要的概念之一，它的引进及其一整套的计算方法，是微积分具有巨大威力的源泉．

直观上说，微商概念来自一个连续变量随另一个连续变量变化的``瞬时''变化率，即函数的变化率，它是质点作变速直线运动的瞬时速度的抽象．

\begin{example}\label{ex:4-1-1}
质点作变速直线运动的瞬时速度．
\end{example}

在第三章 §3 中，我们通过求自由落体瞬时速度引进函数极限概念．现在有了极限概念以后，可以考虑一般的作不等速直线运动的质点的瞬时速度．设质点 $P$ 沿直线作变速运动，用 $s$ 表示从某一选定的时刻开始到时刻 $t$ 为止质点所走过的路程，则 $s$ 是 $t$ 的函数：$s = s(t)$．现在的问题是：已知质点 $P$ 的运动规律 $s = s(t)$，求质点 $P$ 在时刻 $t_0$ 的瞬时速度．

为此，我们考虑 $t_0$ 附近的一段时间间隔——从 $t_0$ 到 $t_0 + \Delta t$，在这段时间内，质点 $P$ 所走过的路程为
\[
\Delta s = s\left(t_0+\Delta t\right) - s\left(t_0\right).
\]
这时平均速度为
\[
\bar{v} = \frac{\Delta s}{\Delta t} = \frac{s\left(t_0+\Delta t\right) - s\left(t_0\right)}{\Delta t}.
\]
由于运动是变速的，像自由落体那样，速度每时每刻都在变化．不过一般说来，当时间间隔 $\left|\Delta t\right|$ 很小时，质点的速度来不及有多大的改变，因此可以把运动近似地看成匀速的．这样，平均速度 $\bar{v} = \dfrac{\Delta s}{\Delta t}$ 就可近似地描述瞬时速度 $v(t_0)$．并且容易理解，当 $\left|\Delta t\right|$ 愈小，则平均速度 $\bar{v} = \dfrac{\Delta s}{\Delta t}$ 就愈接近瞬时速度 $v(t_0)$，因而当 $\Delta t \to 0$ 时，平均速度的极限就是瞬时速度，即
\[
v\left(t_0\right) = \lim_{\Delta t\to 0}\frac{\Delta s}{\Delta t} = \lim_{\Delta t\to 0}\frac{s\left(t_0+\Delta t\right) - s\left(t_0\right)}{\Delta t}.
\]

\begin{example}\label{ex:4-1-2}
非均匀棒的密度．
\end{example}

在物理上，形状接近于直线段的细长物体称为棒．棒的横断面很小而且在任何部位都是一样的．如果棒的任何长度相等的两段都有相同的质量，我们就说这棒是均匀的．对于均匀的棒，它的任何一段的质量，都与它的长度成正比，因此任何一段的质量与它的长度之比都得到同一个常数 $\rho$，这个量 $\rho$ 可以看作是棒的单位长的质量，称为均匀棒的（线）密度．

虽然通常见到的棒的密度大体上都是均匀的，但随着科学技术的发展，如高速流体运动的研究，非均匀密度的概念就越来越重要．为了简单起见，我们以棒为例说明这个概念．

如果一根棒是不均匀的，换句话说，在棒的某些地方物质分布得密一些，在有些地方则不太密，因而棒的同样长度的两段，一般说来就有不同的质量，而每一段的质量与它的长度之比就不相同，很自然地，我们把这个比值称为棒在给定的这一段的平均密度，因为在给定的这一段上，物质的密度还可能有变化，它在这一部分密一些而在另一部分又疏一些，所以知道了这一段的平均密度，一般说来，并不能使我们知道在这一段上的这一点或那一点邻近物质分布得有多密．

我们取棒的一端作为计算的起点，即坐标原点，用 $x$ 表示棒上的点对于计算起点的横坐标．物质分布在 $(0,x)$ 段上的质量是 $x$ 的一个函数，它随 $x$ 的增大而增大；我们把这个函数记作
\[
m = f(x).
\]
取定棒上的一点 $x_0$，给 $x_0$ 一个改变量 $\Delta x$，则分布在由 $x_0$ 到 $x_0 + \Delta x$ 之间的一段上的质量显然是
\[
\Delta m = f\left(x_0+\Delta x\right) - f\left(x_0\right).
\]
因此，棒在这一段上的平均密度是
\[
\bar{\rho} = \frac{\Delta m}{\Delta x} = \frac{f\left(x_0+\Delta x\right) - f\left(x_0\right)}{\Delta x}.
\]
和前面瞬时速度的情况相类似，由于棒是不均匀的，密度在由 $x_0$ 到 $x_0 + \Delta x$ 这一段仍然有变化，所以这个平均密度并不能刻画棒在 $x_0$ 这一点的密度．不过，一般说来，当这段棒的长度 $\left|\Delta x\right|$ 很小时，这个平均密度可以作为棒在 $x_0$ 处密度的近似值．并且容易理解，$\left|\Delta x\right|$ 愈小，则平均密度 $\bar{\rho} = \dfrac{\Delta m}{\Delta x}$ 就愈精确地描写棒在 $x_0$ 的密度情况．因而当 $\Delta x \to 0$ 时，平均密度的极限
\[
\rho\left(x_0\right) = \lim_{\Delta x\to 0}\frac{\Delta m}{\Delta x} = \lim_{\Delta x\to 0}\frac{f\left(x_0+\Delta x\right) - f\left(x_0\right)}{\Delta x}
\]
就完全刻画了棒在 $x_0$ 这一点附近的密度情况，我们称它为棒在 $x_0$ 的线密度．

从上面两个例子可以看出，虽然问题的具体意义不同，一个是用平均速度的极限去刻画瞬时速度，一个是用平均密度的极限去刻画棒在一点的密度，但这两者从数量方面来看，它们都是利用函数 $y = f(x)$ 的改变量与自变量 $x$ 的改变量之比即函数的平均变化速度的极限，来刻画这个函数在一点的变化速度．通常，函数变化的快慢是不均匀的，因此，对一般的函数来说，比式 $\dfrac{\Delta y}{\Delta x}$ 对于不同的 $x_0$ 与 $\Delta x$ 都是不同的，这个比值只刻画了在区间 $[x_0, x_0+\Delta x]$ 或 $[x_0+\Delta x, x_0]$ 上 $y$ 对 $x$ 的平均变化率．对固定的 $\Delta x$ 来说，它只是函数 $y$ 在 $x_0$ 点的变化率的一个近似；而只有在 $\Delta x$ 无限变小的过程中，我们才有可能把这个比值的极限看作函数 $y$ 在 $x_0$ 点的变化率，这个变化率称为微商，把它写成下面的定义．

\begin{definition}\label{def:4-1}
设函数 $y = f(x)$ 在 $x_0$ 点附近有定义．对于自变量在 $x_0$ 点的任一改变量 $\Delta x = x - x_0$，函数在该点的相应改变量为 $\Delta y = f\left(x_0+\Delta x\right) - f\left(x_0\right)$．若极限
\[
\lim_{\Delta x\to 0}\frac{\Delta y}{\Delta x} = \lim_{\Delta x\to 0}\frac{f\left(x_0+\Delta x\right) - f\left(x_0\right)}{\Delta x}
\]
存在，则称函数 $f(x)$ 在 $x_0$ 点可导，并称极限值为 $f(x)$ 在 $x_0$ 点的微商（differential quotient）或导数（derivative），记为
\[
f^{\prime}\left(x_0\right),\quad\text{或}\ \left.y^{\prime}\right|_{x=x_0},\quad\text{或}\ \left.\frac{\mathrm{d}y}{\mathrm{d}x}\right|_{x=x_0}.
\]
若令 $x = x_0 + \Delta x$，则也有
\[
f^{\prime}\left(x_0\right) = \lim_{x\to x_0}\frac{f(x) - f\left(x_0\right)}{x - x_0}.
\]
\end{definition}

函数的微商是一个局部性概念，它描写了在某一点附近函数相对于自变量的变化率，它在数值上的大小反映了函数变化的快慢，而它的符号反映了函数变化的增大或减小的趋势．由于变化速度是一个极其广泛的概念，因此微商的概念在一切自然科学与技术部门都经常遇到．例如电流强度、比热、变力做功的功率、化学反应的速度、有机体的生长速度等．微商和导数这两个名称都带有历史的痕迹．微积分建立初期，没有函数极限的准确概念，只有两个微小量之商这种直观理解，这便是微商这名称的由来．不过，当我们以后有了微分的准确概念后，便知微商乃是微分之商，因此这个名称也有合理之处，这是后话，以后再详细说明．至于导数这个名称，是由于当时``微商''这个概念不精确，受到了责难（这在前面第三章 §1 引进极限，讲自由落体的瞬时速度时提到过），有人把``微商''用其他的方法（代数方法）导出来，所以叫做导数．不过，这种导出的方法并不理想，现在也很少有人用．

回到前面讲的例子，现在可以说，瞬时速度就是路程对时间的微商，线密度就是质量对长度的微商．

\begin{example}\label{ex:4-1-3}
求 $y = x^2$ 在 $x_0$ 的微商．
\end{example}

给自变量以改变量 $\Delta x$，函数有对应的改变量
\[
\Delta y = \left(x_0+\Delta x\right)^2 - x_0^{2} = 2x_0\Delta x + \left(\Delta x\right)^2,
\]
它们之比为 $\dfrac{\Delta y}{\Delta x} = 2x_0 + \Delta x$，令 $\Delta x \to 0$ 取极限，即得函数 $y = x^2$ 在 $x_0$ 的微商
\[
\left.y^{\prime}\right|_{x=x_0} = \lim_{\Delta x\to 0}\frac{\Delta y}{\Delta x} = 2x_0.
\]

\begin{example}\label{ex:4-1-4}
求 $y = \sin x$ 在 $x_0$ 的微商．
\end{example}

当给自变量以改变量 $\Delta x$ 时，函数有改变量
\[
\Delta y = \sin\left(x_0+\Delta x\right) - \sin x_0 = 2\cos\left(x_0+\frac{\Delta x}{2}\right)\sin\frac{\Delta x}{2},
\]
它们之比为
\[
\frac{\Delta y}{\Delta x} = \frac{\cos\left(x_0+\dfrac{\Delta x}{2}\right)\sin\dfrac{\Delta x}{2}}{\dfrac{\Delta x}{2}}.
\]
注意到第三章 §3 讲到的两个重要的极限之一，就有
\[
\lim_{\Delta x\to 0}\frac{\sin\dfrac{\Delta x}{2}}{\dfrac{\Delta x}{2}} = 1,
\]
因此
\[
\lim_{\Delta x\to 0}\frac{\Delta y}{\Delta x} = \cos x_0.
\]
这里用到了 $\cos x$ 在 $x_0$ 的连续性，故 $y = \sin x$ 的微商为
\[
\left.y^{\prime}\right|_{x=x_0} = \cos x_0.
\]

下面讲述微商的几何意义．函数 $y = f(x)$ 在直角坐标系中的图形是一条曲线（图~\ref{fig:4-1}）．设 $y_0 = f(x_0)$，$P(x_0, y_0)$ 是曲线上一点．对应于自变量的改变量 $\Delta x$，函数有改变量 $\Delta y = f(x_0+\Delta x) - f(x_0)$，这样 $Q(x_0+\Delta x, y_0+\Delta y)$ 就是曲线上 $P(x_0, y_0)$ 附近的一点，这时比值 $\dfrac{\Delta y}{\Delta x}$ 就表示了曲线的割线 $PQ$ 的斜率，即 $PQ$ 与 $x$ 轴正方向之间的夹角的正切：
\[
\tan\varphi = \frac{RQ}{PR} = \frac{f\left(x_0+\Delta x\right) - f\left(x_0\right)}{\Delta x}.
\]
当 $\left|\Delta x\right|$ 愈来愈小，$Q$ 点就沿着曲线愈来愈接近于 $P$，而割线 $PQ$ 也就愈来愈接近一条直线．当 $\Delta x \to 0$ 时，$Q$ 就趋向于 $P$，而割线 $PQ$ 的极限位置就是曲线在 $P\left(x_0, y_0\right)$ 的切线，从而夹角 $\varphi$ 也就趋向于切线与 $x$ 轴正向的夹角 $\theta$，割线的斜率便趋向于切线的斜率 $\tan\theta$，即 $f^{\prime}\left(x_0\right) = \tan\theta$．$f^{\prime}\left(x_0\right) > 0$，表示切线与 $x$ 轴正向的夹角为锐角；$f^{\prime}\left(x_0\right) < 0$，表示切线与 $x$ 轴正向的夹角为钝角；$f^{\prime}\left(x_0\right) = 0$，表示切线平行于 $x$ 轴．

\begin{figure}[htbp]
\centering
\includegraphics[width=0.4\textwidth]{fig4-1.png}
\figcap
\label{fig:4-1}
\end{figure}

由此，我们知道曲线 $y = f(x)$ 在 $P(x_0, y_0)$ 处的切线方程为
\[
y - y_0 = f^{\prime}\left(x_0\right)\left(x - x_0\right),
\]
法线方程为
\[
y - y_0 = \frac{-1}{f^{\prime}\left(x_0\right)}\left(x - x_0\right).
\]

\begin{example}\label{ex:4-1-5}
求曲线 $y = x^2$ 在对应于 $x_0 = 2$ 处的切线方程和法线方程．
\end{example}

\begin{proof}[解]
$\left.y^{\prime}\right|_{x=2} = \left.2x\right|_{x=2} = 4$，故在对应于 $x_0 = 2$ 处的切线方程和法线方程分别为
\[
\begin{aligned}
&y - 4 = 4\left(x - 2\right),\ \text{即}\ 4x - y - 4 = 0,\\
&y - 4 = \frac{-1}{4}\left(x - 2\right),\ \text{即}\ x + 4y - 18 = 0.
\end{aligned}
\]
\end{proof}

若 $y = f(x)$ 在一个区间 $(a, b)$ 内每一点都可导，即对区间的每一个 $x_0$，有 $f'(x_0)$ 与之对应，则微商也就构成了一个函数 $y' = f'(x)$，$f'(x_0)$ 就是这个函数在 $x_0$ 点的值．这时称函数 $f(x)$ 在 $(a, b)$ 可导，$f'(x)$ 称为它的导函数，有时也称微商．

\subsection{可导与连续的关系}

从导数的定义可知，若 $f(x)$ 在 $x_0$ 点可导，则极限 $\lim\limits_{\Delta x \to 0} \dfrac{\Delta y}{\Delta x} = f'(x_0)$ 存在，这时 $\dfrac{\Delta y}{\Delta x}$ 与 $f'(x_0)$ 的差是一个无穷小量：
\[
\frac{\Delta y}{\Delta x} = f^{\prime}\left(x_0\right) + \alpha,
\]
其中当 $\Delta x \to 0$ 时，$\alpha \to 0$．因此
\[
\Delta y = f^{\prime}\left(x_0\right)\Delta x + \alpha\Delta x.
\]
故当 $\Delta x \to 0$ 时，有 $\Delta y \to 0$．也就是说 $f(x)$ 在 $x_0$ 点连续．把它写成一个定理：

\begin{theorem}\label{thm:4-1}
若 $f(x)$ 在 $x_0$ 点可导，则 $f(x)$ 在 $x_0$ 点连续．
\end{theorem}

定理说明，函数 $f(x)$ 在 $x_0$ 点不连续，则它在 $x_0$ 点一定不可导，但如果函数在 $x_0$ 点连续，它是否就在 $x_0$ 点可导呢？我们来看下面的反例．

\begin{example}\label{ex:4-1-6}
函数 $f(x) = |x|$，它在 $x = 0$ 点连续，但在 $x = 0$ 点不可导．
\end{example}

\begin{proof}
事实上，
\[
\lim_{\Delta x\to 0^+}\frac{f\left(\Delta x\right) - f(0)}{\Delta x} = \lim_{\Delta x\to 0^+}\frac{\left|0+\Delta x\right| - |0|}{\Delta x} = \lim_{\Delta x\to 0^+}\frac{\Delta x}{\Delta x} = 1,
\]
\[
\lim_{\Delta x\to 0^-}\frac{f\left(\Delta x\right) - f(0)}{\Delta x} = \lim_{\Delta x\to 0^-}\frac{\left|0+\Delta x\right| - |0|}{\Delta x} = \lim_{\Delta x\to 0^-}\frac{-\Delta x}{\Delta x} = -1.
\]
左极限不等于右极限，即差商的极限
\[
\lim_{\Delta x\to 0}\frac{f\left(\Delta x\right) - f(0)}{\Delta x}
\]
不存在，所以 $f(x) = |x|$ 在 $x = 0$ 点不可导．
\end{proof}

\begin{figure}[htbp]
\centering
\includegraphics[width=0.28\textwidth]{fig4-2.png}
\figcap
\label{fig:4-2}
\end{figure}

从图~\ref{fig:4-2} 看到，$y = |x|$ 在 $x = 0$ 是连续的，但它在 $x = 0$ 没有切线，实际上这个点是曲线的尖点．

一般地，若 $\lim\limits_{x\to x_0^-}\dfrac{f(x)-f(x_0)}{x-x_0}$ 存在，则称 $f(x)$ 在 $x_0$ 点有左导数；若 $\lim\limits_{x\to x_0^+}\dfrac{f(x)-f(x_0)}{x-x_0}$ 存在，则称 $f(x)$ 在 $x_0$ 点有右导数．极限值分别记为 $f'_-(x_0)$ 和 $f'_+(x_0)$．显然 $f(x)$ 在 $x_0$ 点可导当且仅当 $f(x)$ 在 $x_0$ 点的左、右导数都存在且相等．

若 $f(x)$ 在区间 $(a, b)$ 可导，而 $f(x)$ 在 $a$ 有右导数 $f'(a)$，在 $b$ 有左导数 $f'(b)$，则称 $f(x)$ 在闭区间 $[a, b]$ 可导．

\subsection{微商的计算}

下面讲述微商的计算．我们要给出基本初等函数的微商公式，并研究函数四则运算、复合函数和反函数的求微商法则．作为一个结果，我们会求一切初等函数的微商（如果它存在的话）．这里特别要注意的是复合函数的微商法则．

$(1)$ 常值函数 $y = c$．

由于 $\Delta y = c - c$ 恒为 $0$，故在任何一点都有
\[
y^{\prime} = \lim_{\Delta x\to 0}\frac{\Delta y}{\Delta x} = \lim_{\Delta x\to 0}\frac{c - c}{\Delta x} = 0,
\]
因此 $\left(c\right)^{\prime} = 0$．

$(2)$ $y = x^n$，其中 $n$ 是正整数．

对任意 $x$，给 $x$ 以改变量 $\Delta x$，对应的函数有改变量
\[
\begin{aligned}
\Delta y &= \left(x+\Delta x\right)^n - x^n\\
&= nx^{n-1}\Delta x + \frac{n(n-1)}{2!}x^{n-2}(\Delta x)^2 + \cdots + \left(\Delta x\right)^n,
\end{aligned}
\]
因此
\[
\lim_{\Delta x\to 0}\frac{\left(x+\Delta x\right)^n - x^n}{\Delta x} = nx^{n-1},
\]
故
\[
\left(x^n\right)^{\prime} = n\,x^{n-1}.
\]

$(3)$ 正弦函数 $y = \sin x$ 与余弦函数 $y = \cos x$．

我们前面的例子已证明过 $\left(\sin x\right)' = \cos x$．同理可推出 $\left(\cos x\right)' = -\sin x$，请读者把证明写出来．

$(4)$ 对数函数 $y = \log_a x$（$a > 0$，$a \neq 1$）．
\[
\begin{aligned}
y^{\prime} &= \lim_{\Delta x\to 0}\frac{\log_a\left(x+\Delta x\right) - \log_a x}{\Delta x} = \lim_{\Delta x\to 0}\frac{\log_a\left(1+\dfrac{\Delta x}{x}\right)}{\dfrac{\Delta x}{x}}\cdot\frac{1}{x}\\
&= \frac{1}{x}\lim_{\Delta x\to 0}\log_a\left(1+\frac{\Delta x}{x}\right)^{\frac{x}{\Delta x}} = \frac{1}{x}\log_a\mathrm{e} = \frac{1}{x\ln a}\;,
\end{aligned}
\]
这里用到了 $\log_a x$ 的连续性以及第三章 §3 讲的两个重要极限之一，$\lim\limits_{x \to 0} (1 + x)^{\frac{1}{x}} = \mathrm{e}$．结果得到
\[
\left(\log_{a}x\right)^{\prime} = \frac{1}{x\ln a}.
\]
特别地
\[
\left(\ln x\right)^{\prime} = \frac{1}{x}.
\]
从 $\sin x$ 与 $\log_a x$ 的微商公式推导过程就知道，我们为什么在第三章 §3 中把那两个极限称为重要极限．

现在我们来证明微商的四则运算法则．

\begin{theorem}\label{thm:4-2}
若函数 $u(x)$ 和 $v(x)$ 在 $x_0$ 点可导，则
\begin{enumerate}[(i)]
\item $\left.\left(u(x) \pm v(x)\right)^{\prime}\right|_{x=x_0} = u'(x_0) \pm v'(x_0)$；
\item $\left.\left(u(x)v(x)\right)^{\prime}\right|_{x=x_0} = u'(x_0)v(x_0) + u(x_0)v'(x_0)$；
\item $\left(\dfrac{u(x)}{v(x)}\right)' = \dfrac{u'(x_0)v(x_0) - u(x_0)v'(x_0)}{v^2(x_0)}$，$\left(v(x_0) \neq 0\right)$．
\end{enumerate}
\end{theorem}

\begin{proof}
$(i)$ 因为
\[
\begin{aligned}
&\frac{u\left(x_0+\Delta x\right) \pm v\left(x_0+\Delta x\right) - \left(u\left(x_0\right) \pm v\left(x_0\right)\right)}{\Delta x}\\
&= \frac{\left[u\left(x_0+\Delta x\right) - u\left(x_0\right)\right] \pm \left[v\left(x_0+\Delta x\right) - v\left(x_0\right)\right]}{\Delta x}\\
&= \frac{\Delta u}{\Delta x} \pm \frac{\Delta v}{\Delta x},
\end{aligned}
\]
所以
\[
\left.\left(u(x) \pm v(x)\right)^{\prime}\right|_{x=x_0} = \lim_{\Delta x \to 0}\left[\frac{\Delta u}{\Delta x} \pm \frac{\Delta v}{\Delta x}\right] = u'(x_0) \pm v'(x_0).
\]
$(ii)$ 由于
\[
\begin{aligned}
&\frac{u\left(x_0+\Delta x\right)v\left(x_0+\Delta x\right) - u\left(x_0\right)v\left(x_0\right)}{\Delta x}\\
&= \frac{u\left(x_0+\Delta x\right) - u\left(x_0\right)}{\Delta x}v\left(x_0+\Delta x\right) + \frac{v\left(x_0+\Delta x\right) - v\left(x_0\right)}{\Delta x}u\left(x_0\right)\\
&= \frac{\Delta u}{\Delta x}v\left(x_0+\Delta x\right) + u\left(x_0\right)\frac{\Delta v}{\Delta x},
\end{aligned}
\]
注意到可导必连续，则
\[
\begin{aligned}
\left.\left(u\left(x\right)v\left(x\right)\right)^{\prime}\right|_{x=x_0} &= \lim_{\Delta x\to 0}\left[\frac{\Delta u}{\Delta x}v\left(x_0+\Delta x\right) + u\left(x_0\right)\frac{\Delta v}{\Delta x}\right]\\
&= u^{\prime}\left(x_0\right)v\left(x_0\right) + u\left(x_0\right)v^{\prime}\left(x_0\right).
\end{aligned}
\]
$(iii)$ 由于
\[
\begin{aligned}
&\frac{\dfrac{u(x_0+\Delta x)}{v(x_0+\Delta x)} - \dfrac{u(x_0)}{v(x_0)}}{\Delta x}\\
&= \frac{u\left(x_0+\Delta x\right)v\left(x_0\right) - u\left(x_0\right)v\left(x_0+\Delta x\right)}{v\left(x_0\right)v\left(x_0+\Delta x\right)\Delta x}\\
&= \frac{\dfrac{u\left(x_0+\Delta x\right) - u\left(x_0\right)}{\Delta x}v\left(x_0\right) - u\left(x_0\right)\dfrac{v\left(x_0+\Delta x\right) - v\left(x_0\right)}{\Delta x}}{v\left(x_0\right)v\left(x_0+\Delta x\right)}\\
&= \frac{\dfrac{\Delta u}{\Delta x}v\left(x_0\right) - u\left(x_0\right)\dfrac{\Delta v}{\Delta x}}{v\left(x_0\right)v\left(x_0+\Delta x\right)},
\end{aligned}
\]
同样利用可导必连续得
\[
\begin{aligned}
\left(\frac{u\left(x\right)}{v\left(x\right)}\right)^{\prime}\bigg|_{x=x_0} &= \lim_{\Delta x\to 0}\frac{\dfrac{\Delta u}{\Delta x}v\left(x_0\right) - u\left(x_0\right)\dfrac{\Delta v}{\Delta x}}{v\left(x_0\right)v\left(x_0+\Delta x\right)}\\
&= \frac{u^{\prime}\left(x_0\right)v\left(x_0\right) - u\left(x_0\right)v^{\prime}\left(x_0\right)}{v^{2}\left(x_0\right)}.
\end{aligned}
\]
定理 4.2 证完．
\end{proof}

利用商的微商运算法则，立得
\[
\begin{aligned}
(\tan x)^{\prime} &= \left(\frac{\sin x}{\cos x}\right)^{\prime} = \frac{(\sin x)^{\prime}\cos x - \sin x(\cos x)^{\prime}}{\cos^{2}x}\\
&= \frac{\cos^{2}x + \sin^{2}x}{\cos^{2}x} = \frac{1}{\cos^{2}x} = \sec^{2}x\;,
\end{aligned}
\]
\[
\begin{aligned}
\left(\cot x\right)^{\prime} &= \left(\frac{\cos x}{\sin x}\right)^{\prime} = \frac{\left(\cos x\right)^{\prime}\sin x - \cos x\left(\sin x\right)^{\prime}}{\sin^{2}x}\\
&= \frac{-\sin^{2}x - \cos^{2}x}{\sin^{2}x} = \frac{-1}{\sin^{2}x} = -\csc^{2}x\;.
\end{aligned}
\]

最后我们将利用反函数的求导法则来求指数函数和反三角函数的微商．先给出反函数的求微商法则．

\begin{theorem}\label{thm:4-3}
若函数 $y = f(x)$ 在 $x_0$ 点附近连续且严格单调，又 $f'(x_0) \neq 0$，则其反函数 $x = \varphi(y)$ 在点 $y_0 = f(x_0)$ 可导，且
\[
\varphi^{\prime}\left(y_0\right) = \frac{1}{f^{\prime}\left(x_0\right)}.
\]
\end{theorem}

\begin{proof}
由 $f(x)$ 在 $x_0$ 附近连续且严格单调，则反函数 $x = \varphi(y)$ 在 $y_0$ 点附近连续且严格单调．因此，若 $y - y_0 \neq 0$，则 $x - x_0 = \varphi(y) - \varphi(y_0) \neq 0$，且当 $y \to y_0$ 时有 $x \to x_0$，故由复合函数求极限法则得
\[
\begin{aligned}
\lim_{y\to y_0}\frac{\varphi(y) - \varphi\left(y_0\right)}{y - y_0} &= \lim_{y\to y_0}\frac{\varphi(y) - \varphi\left(y_0\right)}{f(\varphi(y)) - f\left(\varphi\left(y_0\right)\right)}\\
&= \lim_{x\to x_0}\frac{x - x_0}{f(x) - f\left(x_0\right)} = \lim_{x\to x_0}\frac{1}{\dfrac{f(x) - f\left(x_0\right)}{x - x_0}} = \frac{1}{f^{\prime}\left(x_0\right)}.
\end{aligned}
\]
定理 4.3 证完．
\end{proof}

$(5)$ 指数函数 $y = a^x$（$a > 0$，$a \neq 1$）．

指数函数 $y = a^x$，$x \in (-\infty, +\infty)$ 是对数函数 $x = \log_a y$，$y \in (0, +\infty)$ 的反函数，而 $x = \log_a y$ 在定义域内连续、严格单调且 $x' = \dfrac{1}{y\ln a} \neq 0$，因此 $y = a^x$ 在其定义域内可导，且
\[
y^{\prime} = \frac{1}{x^{\prime}\left(y\right)} = y\ln a = a^{x}\ln a.
\]

$(6)$ 反三角函数．

\begin{enumerate}[(i)]
\item 若 $y = \arcsin x$，则 $y' = \dfrac{1}{\sqrt{1 - x^2}}\ \left(-\dfrac{\pi}{2} < y < \dfrac{\pi}{2},\ -1 < x < 1\right)$；
\item 若 $y = \arccos x$，则 $y' = \dfrac{-1}{\sqrt{1 - x^2}}\ \left(0 < y < \pi,\ -1 < x < 1\right)$；
\item 若 $y = \arctan x$，则 $y' = \dfrac{1}{1 + x^2}\ \left(-\dfrac{\pi}{2} < y < \dfrac{\pi}{2},\ -\infty < x < +\infty\right)$；
\item 若 $y = \operatorname{arccot} x$，则 $y' = \dfrac{-1}{1 + x^2}\ \left(0 < y < \pi,\ -\infty < x < +\infty\right)$．
\end{enumerate}

事实上，由于 $y = \arcsin x$，$x \in (-1,1)$ 是 $x = \sin y$，$y \in \left(-\dfrac{\pi}{2}, \dfrac{\pi}{2}\right)$ 的反函数，而 $x = \sin y$ 在 $\left(-\dfrac{\pi}{2}, \dfrac{\pi}{2}\right)$ 连续、严格单调且 $x' = \cos y \neq 0$，因此
\[
y^{\prime} = \frac{1}{x^{\prime}(y)} = \frac{1}{\cos y} = \frac{1}{\sqrt{1 - \sin^{2}y}} = \frac{1}{\sqrt{1 - x^{2}}},
\]
故 (i) 成立．(ii) 同理可证．

由于 $y = \arctan x$，$x \in (-\infty, +\infty)$ 是 $x = \tan y$，$y \in \left(-\dfrac{\pi}{2}, \dfrac{\pi}{2}\right)$ 的反函数，而 $x = \tan y$ 在 $\left(-\dfrac{\pi}{2}, \dfrac{\pi}{2}\right)$ 连续、严格单调且 $x' = \sec^2 y \neq 0$，因此
\[
y^{\prime} = \frac{1}{x^{\prime}(y)} = \frac{1}{\sec^{2}y} = \frac{1}{1 + \tan^{2}y} = \frac{1}{1 + x^{2}},
\]
故 (iii) 成立．(iv) 同理可证．

下面叙述复合函数的求微商法则．

\begin{theorem}[复合函数求导法则]\label{thm:4-4}
若函数 $u = g(x)$ 在 $x_0$ 点可导，$y = f(u)$ 在 $u_0 = g(x_0)$ 点可导，则复合函数 $f \circ g(x)$ 在 $x_0$ 点可导，且
\[
\left(f\circ g\right)^{\prime}\left(x_0\right) = f^{\prime}\left(u_0\right)g^{\prime}\left(x_0\right) = f^{\prime}\left(g\left(x_0\right)\right)g^{\prime}\left(x_0\right),
\]
或
\[
\frac{\mathrm{d}y}{\mathrm{d}x}\bigg|_{x=x_0} = \frac{\mathrm{d}y}{\mathrm{d}u}\bigg|_{u=u_0}\cdot\frac{\mathrm{d}u}{\mathrm{d}x}\bigg|_{x=x_0}.
\]
\end{theorem}

\begin{proof}
定义函数
\[
H\left(u\right) = \begin{cases}
\dfrac{f\left(u\right) - f\left(u_0\right)}{u - u_0}, & u \neq u_0,\\[2ex]
f^{\prime}\left(u_0\right), & u = u_0,
\end{cases}
\]
则
\[
\lim_{u\to u_0}H(u) = \lim_{u\to u_0}\frac{f(u) - f\left(u_0\right)}{u - u_0} = f^{\prime}\left(u_0\right) = H\left(u_0\right),
\]
故 $H(u)$ 在 $u_0$ 点连续．在恒等式
\[
f(u) - f\left(u_0\right) = H(u)\left(u - u_0\right)
\]
中令 $u = g(x)$ 代入得
\[
f(g(x)) - f\left(g\left(x_0\right)\right) = H(g(x))\left(g(x) - g\left(x_0\right)\right),
\]
两边除以 $x - x_0$ 得
\[
\frac{f(g(x)) - f\left(g\left(x_0\right)\right)}{x - x_0} = H(g(x))\frac{g(x) - g\left(x_0\right)}{x - x_0}.
\]
由复合函数的连续性知
\[
\lim_{x\to x_0}H(g(x)) = H\left(g\left(x_0\right)\right) = H\left(u_0\right) = f^{\prime}\left(u_0\right),
\]
由于 $g(x)$ 在 $x_0$ 点可导
\[
\lim_{x\to x_0}\frac{g(x) - g\left(x_0\right)}{x - x_0} = g^{\prime}\left(x_0\right),
\]
故
\[
\begin{aligned}
\left(f^{\circ}g\right)^{\prime}\left(x_0\right) &= \lim_{x\to x_0}\frac{f\left(g\left(x\right)\right) - f\left(g\left(x_0\right)\right)}{x - x_0}\\
&= f^{\prime}\left(u_0\right)g^{\prime}\left(x_0\right) = f^{\prime}\left(g\left(x_0\right)\right)g^{\prime}\left(x_0\right).\qedhere
\end{aligned}
\]
\end{proof}

请读者思考：为什么要引入函数 $H(u)$？下面的证明是否正确：
\[
\begin{aligned}
\frac{f(g(x)) - f\left(g\left(x_0\right)\right)}{x - x_0} &= \frac{f(u) - f\left(u_0\right)}{u - u_0}\cdot\frac{u - u_0}{x - x_0}\\
&= \frac{f(u) - f\left(u_0\right)}{u - u_0}\cdot\frac{g(x) - g\left(x_0\right)}{x - x_0},
\end{aligned}
\]
当 $x \to x_0$ 时，有 $u \to u_0$，故定理结论成立．

若 $y = f(u)$ 的定义域包含 $u = g(x)$ 的值域，且两个函数在各自的定义域上可导，则复合函数 $f \circ g(x)$ 在定义域上可导，导函数为
\[
\left(f\circ g\right)^{\prime}\left(x\right) = f^{\prime}\left(u\right)g^{\prime}\left(x\right) = f^{\prime}\left(g\left(x\right)\right)g^{\prime}\left(x\right),
\]
或
\[
\frac{\mathrm{d}y}{\mathrm{d}x} = \frac{\mathrm{d}y}{\mathrm{d}u}\frac{\mathrm{d}u}{\mathrm{d}x}.
\]
也可写成
\[
y^{\prime}{}_{x} = y^{\prime}{}_{u}\cdot u_{x}^{\prime}.
\]

上述公式也称为链式法则，它可推广到多个函数的复合情形．例如，可导函数 $y = f(u)$，$u = g(v)$ 和 $v = h(x)$ 的复合函数 $y = (f \circ g \circ h)(x)$ 对 $x$ 的导数为
\[
\left(f\circ g\circ h\right)^{\prime}(x) = f^{\prime}(u)g^{\prime}(v)h^{\prime}(x) = f^{\prime}(g(h(x)))g^{\prime}(h(x))h^{\prime}(x),
\]
或
\[
y^{\prime}{}_{x} = y_{u}{}^{\prime}u_{v}{}^{\prime}v_{x}{}^{\prime}.
\]

链式法则有很简单的物理意义．$g'(x)$ 是 $x$ 在单位时间内 $u = g(x)$ 的变化，$f'(u)$ 是 $u$ 在单位时间内 $y = f(u)$ 的变化，而 $(f \circ g)'(x)$ 是 $x$ 在单位时间由 $y = f(g(x))$ 决定的变化，所以应有
\[
\left(f\circ g\right)^{\prime}(x) = f^{\prime}(u)g^{\prime}(x).
\]

$(7)$ 幂函数 $y = x^{\alpha}$ 在 $x > 0$ 的微商 $y' = \alpha x^{\alpha-1}$．

事实上 $y = x^{\alpha} = \mathrm{e}^{\alpha\ln x}$，它可视为 $y = \mathrm{e}^{u}$ 和 $u = \alpha\ln x$ 的复合．由链式法则有
\[
\left(x^{\alpha}\right)^{\prime} = \left(\mathrm{e}^{u}\right)^{\prime}\left(\alpha\ln x\right)^{\prime} = \mathrm{e}^{u}\frac{\alpha}{x} = \mathrm{e}^{\alpha\ln x}\frac{\alpha}{x} = \alpha x^{\alpha-1}.
\]

现在把求微商的方法总结一下．首先，我们有一张微商表，它包含了所有基本初等函数的微商公式：
\begin{enumerate}[(1)]
\item $\left(c\right)^{\prime} = 0$；
\item $\left(x^{\alpha}\right)^{\prime} = \alpha x^{\alpha-1}$，$\left(x^n\right)^{\prime} = nx^{n-1}$，
\[
\left(\frac{1}{x}\right)^{\prime} = -\frac{1}{x^{2}},\qquad\left(\sqrt{x}\right)^{\prime} = \frac{1}{2\sqrt{x}};
\]
\item $\left(a^x\right)^{\prime} = a^x\ln a$，$\left(\mathrm{e}^x\right)^{\prime} = \mathrm{e}^x$；
\item $\left(\log_{a}x\right)^{\prime} = \dfrac{1}{x\ln a}$，$\left(\ln x\right)^{\prime} = \dfrac{1}{x}$；
\item $\left(\sin x\right)' = \cos x$；
\item $\left(\cos x\right)' = -\sin x$；
\item $\left(\tan x\right)' = \dfrac{1}{\cos^2 x} = \sec^2 x$；
\item $\left(\cot x\right)' = \dfrac{-1}{\sin^2 x} = -\csc^2 x$；
\item $\left(\arcsin x\right)' = \dfrac{1}{\sqrt{1 - x^2}}$；
\item $\left(\arccos x\right)' = -\dfrac{1}{\sqrt{1 - x^2}}$；
\item $\left(\arctan x\right)' = \dfrac{1}{1 + x^2}$；
\item $\left(\operatorname{arccot} x\right)' = -\dfrac{1}{1 + x^2}$．
\end{enumerate}

其次，我们有微商的运算法则：
\begin{enumerate}[(1)]
\item $\left(u \pm v\right)' = u' \pm v'$；
\item $\left(uv\right)' = u'v + uv'$；
\item $\left(\dfrac{u}{v}\right)^{\prime} = \dfrac{u^{\prime}v - uv^{\prime}}{v^{2}}$；
\item $y_{x}^{\prime} = \dfrac{1}{x_{y}^{\prime}}$；
\item $y'_{x} = y'_{u} u'_{x}$．
\end{enumerate}

这张微商表和这些求微商法则，读者是要熟记的．根据这张微商表和微商法则，我们便能够计算一切初等函数的微商了．由此可以体会到作为函数``瞬时''变化率的微商概念的重要性．设想如果人们只有平均变化率，也就是只有差商 $\dfrac{\Delta y}{\Delta x}$，而不令 $\Delta x \to 0$ 取极限，人们可以得到像 $\left(\sin x\right)' = \cos x$，$\left(\mathrm{e}^{\tan x}\right)' = \dfrac{1}{\cos^2 x}\mathrm{e}^{\tan x}$ 这样简单的计算公式吗？可以这样说，能够计算，或更准确地，能够``机械化''地计算一切初等函数的微商，是引入微商概念（瞬时变化率）的重要原因，这一点等学到本书第七章、第八章时，读者将会更深刻地领会这一点．

\begin{example}\label{ex:4-1-7}
设 $y = \sin x^2$，求 $y'$．
\end{example}

\begin{proof}[解]
$y = \sin x^2$ 可视为 $y = \sin u$ 和 $u = x^2$ 的复合，故
\[
y^{\prime} = \left(\sin u\right)^{\prime}\left(x^{2}\right)^{\prime} = (\cos u)\cdot 2x = 2x\cos x^{2}.
\]
\end{proof}

\begin{example}\label{ex:4-1-8}
设 $y = \sin^2 x$，求 $y'$．
\end{example}

\begin{proof}[解]
$y = \sin^2 x$ 可视为 $y = u^2$ 和 $u = \sin x$ 的复合，故
\[
y^{\prime} = \left(u^{2}\right)^{\prime}(\sin x)^{\prime} = 2u\cos x = 2\sin x\cos x = \sin 2x.
\]
\end{proof}

\begin{example}\label{ex:4-1-9}
设 $y = \dfrac{1}{\sqrt{1 - x^{2}}}$，求 $y'$．
\end{example}

\begin{proof}[解]
$y = \dfrac{1}{\sqrt{1 - x^{2}}} = \left(1 - x^{2}\right)^{-\frac{1}{2}}$ 可视为 $y = u^{-\frac{1}{2}}$ 和 $u = 1 - x^2$ 的复合，故
\[
y^{\prime} = \left(u^{-\frac{1}{2}}\right)^{\prime}\left(1 - x^{2}\right)^{\prime} = -\frac{1}{2}u^{-\frac{3}{2}}\left(-2x\right) = \frac{x}{\left(1 - x^{2}\right)^{\frac{3}{2}}}.
\]
\end{proof}

\begin{example}\label{ex:4-1-10}
设 $y = \dfrac{x + \sqrt{1 + x^{2}}}{(x + 1)^{2} \sqrt{1 + x^{2}}}\ (x > -1)$，求 $y'$．
\end{example}

\begin{proof}[解]
两边取对数得
\[
\begin{aligned}
\ln y &= \ln\left(x+\sqrt{1+x^{2}}\right) - \ln\left(x+1\right)^{2} - \ln\sqrt{1+x^{2}}\\
&= \ln\left(x+\sqrt{1+x^{2}}\right) - 2\ln\left(x+1\right) - \frac{1}{2}\ln\left(1+x^{2}\right),
\end{aligned}
\]
上式两边对 $x$ 求导得
\[
\begin{aligned}
\frac{y^{\prime}}{y} &= \frac{1}{x+\sqrt{1+x^{2}}}\left(1+\frac{2x}{2\sqrt{1+x^{2}}}\right) - \frac{2}{x+1} - \frac{1}{2}\frac{2x}{1+x^{2}}\\
&= \frac{1}{\sqrt{1+x^{2}}} - \frac{2}{x+1} - \frac{x}{1+x^{2}},
\end{aligned}
\]
因此
\[
y^{\prime} = \frac{x+\sqrt{1+x^{2}}}{(x+1)^{2}\sqrt{1+x^{2}}}\left(\frac{1}{\sqrt{1+x^{2}}} - \frac{2}{x+1} - \frac{x}{1+x^{2}}\right).
\]
\end{proof}

\begin{example}\label{ex:4-1-11}
设 $y = x^{\sin x}\ (x > 0)$，求 $y'$．
\end{example}

\begin{proof}[解]
两边取对数
\[
\ln y = \sin x\ln x,
\]
再两边对 $x$ 求导得
\[
\frac{y^{\prime}}{y} = \cos x\cdot\ln x + \frac{\sin x}{x},
\]
故
\[
y^{\prime} = x^{\sin x}\left(\cos x\cdot\ln x + \frac{\sin x}{x}\right).
\]
\end{proof}

例 10 和例 11 采用的方法也称为对数求导法，它能简化求导运算．例 11 也可用链式法则求得．因为 $y = x^{\sin x} = \mathrm{e}^{\sin x\ln x}$，所以
\[
\begin{aligned}
y^{\prime} &= \mathrm{e}^{\sin x\cdot\ln x}\left(\sin x\cdot\ln x\right)^{\prime}\\
&= \mathrm{e}^{\sin x\cdot\ln x}\left(\cos x\cdot\ln x + \frac{\sin x}{x}\right)\\
&= x^{\sin x}\left(\cos x\cdot\ln x + \frac{\sin x}{x}\right).
\end{aligned}
\]

从这几个例子，读者可以体会到，我们能够计算一切初等函数的微商了，而且计算不需要任何技巧，只用到微商表与运算法则，的确是``机械化''了．

下面再举两个说明函数在一点连续但并不可导的例子．

\begin{example}\label{ex:4-1-12}
函数 $f(x) = (x-1)^{\frac{1}{3}}$，$x\in(-\infty,+\infty)$ 是初等函数，故在定义域内连续，但
\[
\lim_{x\to 1}\frac{(x-1)^{\frac{1}{3}}}{x-1} = \lim_{x\to 1}\left(x-1\right)^{-\frac{2}{3}} = \infty,
\]
故 $f(x)$ 在 $x = 1$ 点不可导．当 $x \neq 1$ 时有
\[
f^{\prime}\left(x\right) = \frac{1}{3}\left(x-1\right)^{-\frac{2}{3}}.
\]
几何上表示曲线在 $x = 1$ 处的切线平行于 $y$ 轴．
\end{example}

\begin{example}\label{ex:4-1-13}
设
\[
f(x) = \begin{cases}
x\sin\dfrac{1}{x}, & x \neq 0,\\
0, & x = 0.
\end{cases}
\]
\end{example}

当 $x \neq 0$ 时，函数 $f(x)$ 是可导的：
\[
f^{\prime}\left(x\right) = \sin\frac{1}{x} - \frac{1}{x}\cos\frac{1}{x}.
\]
显然 $f(x)$ 在 $x = 0$ 连续．而由于极限
\[
\lim_{x\to 0}\frac{f(x) - f(0)}{x} = \lim_{x\to 0}\frac{x\sin\dfrac{1}{x}}{x} = \lim_{x\to 0}\sin\frac{1}{x}
\]
不存在，故 $f(x)$ 在 $x = 0$ 点不可导．我们知道，当 $x \to 0$ 时，$\sin\dfrac{1}{x}$ 不断地在 $1$ 和 $-1$ 之间摆动．从图形上看就是当 $Q$ 点沿曲线趋于原点时，割线 $OQ$ 在直线 $y = \pm x$ 之间摆动．

\begin{figure}[htbp]
\centering
\includegraphics[width=0.42\textwidth]{fig4-3.png}
\figcap
\label{fig:4-3}
\end{figure}

注意，并不是割线不断摆动就无切线．例如函数
\[
f(x) = \begin{cases}
x^{2}\sin\dfrac{1}{x}, & x \neq 0,\\
0, & x = 0,
\end{cases}
\]
有
\[
\lim_{x\to 0}\frac{f(x) - f(0)}{x} = \lim_{x\to 0}x\sin\frac{1}{x} = 0,
\]
故
\[
f^{\prime}\left(x\right) = \begin{cases}
2x\sin\dfrac{1}{x} - \cos\dfrac{1}{x}, & x \neq 0,\\
0, & x = 0.
\end{cases}
\]
可见 $f(x)$ 在 $x = 0$ 点可导，事实上在 $0$ 点割线的斜率 $x\sin\dfrac{1}{x}$ 也是不断摆动的，但它有个极限位置 $y = 0$．

上面几个例子都说明，函数可以在一点（当然也可以在任意有穷个点）连续但不可导，这在几何上是很直观的．是否存在函数，在整个数轴上连续，但处处都不可导呢？历史上很长一段时间内，人们从几何直观出发，总认为这是不可能的．怎么能想像一条连续曲线处处都没有切线？直至 $19$ 世纪中叶，魏尔斯特拉斯给出了一个例子，它在数轴上是处处连续处处不可导的，才使人们纠正了过去的误解．这例子表明，单靠直观，人们的认识是会受到局限的．

\begin{figure}[htbp]
\centering
\includegraphics[width=0.5\textwidth]{fig4-4.png}
\figcap
\label{fig:4-4}
\end{figure}

\begin{xiti}
\item 求抛物线 $y = x^2$ 在 $A(1, 1)$ 点和在 $B(-2, 4)$ 点的切线方程和法线方程．

\item 若 $s = vt - \frac{1}{2}gt^2$，求：
\begin{ziti}
\item 在 $t = 1$ 与 $t = 1 + \Delta t$ 之间的平均速度（设 $\Delta t = 1, 0.1, 0.01$）；
\item 在 $t = 1$ 的瞬时速度．
\end{ziti}

\item 试确定曲线 $y = \ln x$ 在哪些点的切线平行于下列直线：
\begin{ziti}
\item $y = x - 1$；
\item $y = 2x - 3$．
\end{ziti}

\item 设 $f(x) = \begin{cases} x^2, & x \geqslant 3, \\ ax + b, & x < 3, \end{cases}$ 试确定 $a, b$ 的值，使 $f(x)$ 在 $x = 3$ 处可导．

\item 求下列曲线在指定点 $P$ 的切线方程和法线方程：
\begin{ziti}
\item $y = \frac{x^2}{4}$，$P(2, 1)$；
\item $y = \cos x$，$P(0, 1)$．
\end{ziti}

\item 求下列函数的导函数：
\begin{ziti}
\item $f(x) = |x|^3$；
\item $f(x) = \begin{cases} x + 1, & x \geqslant 0, \\ 1, & x < 0 \end{cases}$．
\end{ziti}

\item 设函数 $f(x) = \begin{cases} x^m \sin\frac{1}{x}, & x \neq 0, \\ 0, & x = 0 \end{cases}$（$m$ 为正整数）．试问：
\begin{ziti}
\item $m$ 等于何值时，$f(x)$ 在 $x = 0$ 连续；
\item $m$ 等于何值时，$f(x)$ 在 $x = 0$ 可导；
\item $m$ 等于何值时，$f'(x)$ 在 $x = 0$ 连续．
\end{ziti}

\item 设 $g(0) = g'(0) = 0$，$f(x) = \begin{cases} g(x)\sin\frac{1}{x}, & x \neq 0, \\ 0, & x = 0, \end{cases}$ 求 $f'(0)$．

\item 证明：若 $f'(x_0)$ 存在，则
\[
\lim_{\Delta x \to 0} \frac{f(x_0 + \Delta x) - f(x_0 - \Delta x)}{2\Delta x} = f'(x_0).
\]

\item 设 $f(x)$ 是定义在 $(-\infty, +\infty)$ 上的函数，且对任意 $x_1, x_2 \in (-\infty, +\infty)$，有
\[
f(x_1 + x_2) = f(x_1)f(x_2).
\]
若 $f'(0) = 1$，证明对任意 $x \in (-\infty, +\infty)$，有 $f'(x) = f(x)$．

\item 设 $f(x)$ 是偶函数，且 $f'(0)$ 存在，证明：$f'(0) = 0$．

\item 设 $f(x)$ 是奇函数，且 $f'(x_0) = 3$，求 $f'(-x_0)$．

\item 用定义证明：可导的偶函数的导函数是奇函数，而可导的奇函数的导函数是偶函数．

\item 求下列函数的导数：
\begin{ziti}
\item $y = x^2 \sin x$；
\item $y = x \cos x + 3x^2$；
\item $y = x \tan x - 7x + 6$；
\item $y = e^x \sin x - 7\cos x + 5x^2$；
\item $y = 4\sqrt{x} + \frac{1}{x} - 2x^3$；
\item $y = 3x + 5\sqrt{x} + \frac{7}{x^3}$；
\item $y = \frac{1 + x^2}{1 - x^2}$；
\item $y = \frac{1}{1 + x + x^2}$；
\item $y = \frac{x}{(1-x)(2-x)}$；
\item $y = \frac{1}{1 + \sqrt{x}} - \frac{1}{1 - \sqrt{x}}$；
\item $y = \frac{1 + \sqrt{x}}{1 - \sqrt{x}}$；
\item $y = \frac{1}{3\sqrt{x}} + \sqrt[3]{x}$；
\item $y = x^3 \ln x - \frac{1}{n}x^n$；
\item $y = \frac{\cos x}{x^4}\ln\frac{1}{x}$；
\item $y = \left(x + \frac{1}{x}\right)\ln x$；
\item $y = \frac{x\cos x - \ln x}{x + 1}$；
\item $y = \frac{1}{x + \cos x}$；
\item $y = \frac{x\sin x + \cos x}{x\sin x - \cos x}$；
\item $y = \frac{xe^x - 1}{\sin x}$；
\item $y = x\sin x \ln x$．
\end{ziti}

\item 求下列复合函数的导数：
\begin{ziti}
\item $y = (x^3 - 4)^3$；
\item $y = x(a^2 + x^2)\sqrt{a^2 - x^2}$；
\item $y = \frac{x}{\sqrt{a^2 - x^2}}$；
\item $y = \sqrt[3]{\frac{1 + x^3}{1 - x^3}}$；
\item $y = \ln(\ln x)$；
\item $y = \frac{1}{2}\ln\left|\frac{a+x}{a-x}\right|$；
\item $y = \ln\left(x + \sqrt{a^2 + x^2}\right)$；
\item $y = \ln\tan\frac{x}{2}$；
\item $y = \cos(\cos\sqrt{x})$；
\item $y = \cos^3 x - \cos 3x$；
\item $y = \frac{1}{\sqrt{2\pi}}e^{-3x^2}$；
\item $y = \arcsin(\sin x \cos x)$；
\item $y = \arctan\frac{2x}{1-x^2}$；
\item $y = e^{-x^2 + 2x}$；
\item $y = \ln\sqrt{\frac{(x+2)(x+3)}{x+1}}$；
\item $y = e^{2x}\sin 3x + \frac{x^2}{2}$；
\item $y = \frac{e^{-kx}\sin \omega x}{1+x}$（$k, \omega$ 为常数）；
\item $y = x\sqrt{a^2 - x^2} + \frac{x}{\sqrt{a^2 - x^2}}$；
\item $y = \sin^n x \cos nx$；
\item $y = \ln\frac{\sqrt{1+x} - \sqrt{1-x}}{\sqrt{1+x} + \sqrt{1-x}}$．
\end{ziti}

\item 用对数求导法求下列函数的导数：
\begin{ziti}
\item $y = x\sqrt{\frac{1-x}{1+x}}$；
\item $y = \frac{x^2}{1-x}\sqrt{\frac{1+x}{1+x+x^2}}$；
\item $y = (x + \sqrt{1+x^2})^n$；
\item $y = x^x$ ($x > 0$)；
\item $y = x^{\ln x}$ ($x > 0$)；
\item $y = (1+x)^{\frac{1}{x}}$ ($x > 0$)；
\item $y = x^{\tan x}$ ($x > 0$)；
\item $y = a^{\sin x}$ ($a > 0$)．
\end{ziti}

\item 设 $f(x)$ 是对 $x$ 可导的函数，求 $\frac{\dd y}{\dd x}$：
\begin{ziti}
\item $y = f(x^2)$；
\item $y = f(e^x)e^{f(x)}$；
\item $y = f(f(f(x)))$．
\end{ziti}

\item 设 $\varphi(x)$ 和 $\psi(x)$ 是对 $x$ 可导的函数，求 $\frac{\dd y}{\dd x}$：
\begin{ziti}
\item $y = \sqrt{\varphi^2(x) + \psi^2(x)}$；
\item $y = \arctan\frac{\varphi(x)}{\psi(x)}$ ($\psi(x) \neq 0$)；
\item $y = \sqrt[4]{x}\sqrt{\psi(x)}$ ($\varphi(x) \neq 0, \psi(x) > 0$)；
\item $y = \log_{\varphi(x)}\psi(x)$ ($\varphi(x) > 0, \psi(x) > 0, \varphi(x) \neq 1$)．
\end{ziti}

\item 求下列函数的导数：
\begin{ziti}
\item $y = e^{ax}(\cos bx + \sin bx)$；
\item $y = x\arctan x - \frac{1}{2}\ln(1+x^2)$；
\item $y = \arctan\frac{\sqrt{1-x^2}-1}{x} + \arctan\frac{2x}{1-x^2}$；
\item $y = \arctan(\tan^2 x)$；
\item $y = \left(\frac{a}{b}\right)^x \left(\frac{b}{x}\right)^a \left(\frac{x}{a}\right)^b$ ($a, b > 0$)；
\item $y = \frac{x}{2}\sqrt{a^2-x^2} + \frac{a^2}{2}\arcsin\frac{x}{a}$ ($a > 0$)；
\item $y = \frac{x}{2}\sqrt{a^2+x^2} + \frac{a^2}{2}\ln|x + \sqrt{a^2+x^2}|$ ($a > 0$)；
\item $y = \ln\left(\arccos\frac{1}{\sqrt{x}}\right)$；
\item $y = x^{a^a} + a^{x^a} + a^{a^x}$ ($a > 0$)；
\item $y = \frac{1}{6}\ln\frac{(x+1)^2}{x^2-x+1} + \frac{1}{\sqrt{3}}\arctan\frac{2x-1}{\sqrt{3}}$．
\end{ziti}
\end{xiti}
